\documentclass[10pt,a4paper]{amsart}
\usepackage[margin=19mm]{geometry}
\usepackage{amsmath,amssymb,mathtools}
\usepackage[scr=rsfs,cal=euler]{mathalfa}
\usepackage{xfrac}
\usepackage[unicode,hidelinks,bookmarks=false]{hyperref}
\newcommand{\ie}{i.e.\ }
\newcommand{\cf}{cf.\ }
\newcommand{\resp}{respectively\ }
\newcommand{\Subsection}{\subsection}
\providecommand{\nobreakdash}{\nobreak\hbox{-}}
\setlength{\emergencystretch}{2em}
\multlinegap0pt
\usepackage{mathtools}
\usepackage{xcolor}
\usepackage{tikz}
\usepackage{amssymb}
\usepackage{braket}
\newtheorem{prop}{Proposition}
\DeclareMathOperator{\id}{id}
\newcommand{\sh}[1]{{\ensuremath{\shspace\hspace{1mm}\makebox[-1mm]{$\langle$}\makebox[0mm]{$\langle$}\hspace{1mm}{#1}\makebox[1mm]{$\rangle$}\makebox[0mm]{$\rangle$}\shspace}}}
\newcommand{\shspace}{\hspace{.5mm}} %to make it easier to read when there are many shadows/dual shadows together
\DeclareMathOperator{\tr}{tr}
\title{Twisted Bicategorical Shadows and Traces}
\author{Zhonghui Sun}
\date{Source: arXiv:2609.04660v1 (CC BY 4.0)}
\begin{document}
\maketitle

\noindent\emph{Source and licence.} Twisted Bicategorical Shadows and Traces. Version arXiv:2609.04660v1 (CC BY 4.0), distributed by arXiv under the Creative Commons Attribution 4.0 International licence, http://creativecommons.org/licenses/by/4.0/, as recorded in the arXiv licence metadata for this exact version and shown on the abstract page https://arxiv.org/abs/2609.04660v1. Full licence notice: \texttt{licenses/arxiv-2609.04660-CC-BY-4.0.txt}.\par\smallskip

\noindent\textit{Editorial scope.} Counted unit: the article's prop prop:tw-multiplicativity (source0.tex lines 5344--5355). The article's complete original proof of it is delivered with it (source0.tex lines 5357--5461, one \texttt{proof} environment of 105 lines). No mathematics is written, repaired or re-derived: the statement and the proof are the article's own lines, in the article's own notation, and no retained line is substituted or re-flowed.  No further source-local context is needed: every cross-reference of every retained line resolves inside the two delivered spans.  Nothing was omitted from any retained span, so the package carries no omission record.  No retained line contains a citation, so the wrapper carries no bibliography and no bibliography entry.  No retained line contains a figure environment, an image-inclusion command or drawing code, so no image is delivered and the package declares no asset dependency.  Editorial formatting only, in addition: an AMS \texttt{amsart} preamble that restates the article's own declarations and macros used by the retained lines, namely the environments \texttt{prop} on separate counters, together with the standard packages those lines need.  The retained environments print in this document as Proposition 1 in order of appearance, whereas the article prints the same items with the numbering of its own full document.  The complete native source of the article as distributed in the arXiv e-print for this exact version is bundled unchanged as \texttt{sources/209380/source0.tex}.
\begin{prop}[Twisted multiplicativity]\label{prop:tw-multiplicativity}
Let
\(X\in\mathcal B^{G\text{-tw}}(A,B)\) and
\(Y\in\mathcal B^{G\text{-tw}}(B,C)\) be left dualizable, and let
\(f_X\colon X\Rightarrow X\) and \(f_Y\colon Y\Rightarrow Y\).  Then
\(X\odot Y\) is left dualizable and
\[
\tr^g(f_X\odot f_Y)
=\tr^g(f_X)\circ\tr^g(f_Y)
\colon{}^g\!\sh{U_C}\longrightarrow{}^g\!\sh{U_A}.
\]
\end{prop}

\begin{proof}
Choose left duals \({}^*X\) and \({}^*Y\).  Then \(X\odot Y\)
has left dual \({}^*Y\odot{}^*X\), with
\[
\bar\eta_{X\odot Y}
=
(\id_{{}^*Y}\odot\bar\eta_X\odot\id_Y)
\circ\bar\eta_Y
\quad \text{and} \quad
\bar\epsilon_{X\odot Y}
=
\bar\epsilon_X\circ
(\id_X\odot\bar\epsilon_Y\odot\id_{{}^*X}).
\]
Suppressing associativity and unit constraints, put
\(
a_X
:=
(\id_{{}^*X}\odot f_X)\circ\bar\eta_X
\colon
U_B\Rightarrow{}^*X\odot X.
\)
After substituting the above duality maps into the definition of the
twisted trace, we obtain
\begin{align*}
\tr^g(f_X\odot f_Y)
={}&
{}^g\!\sh{\bar\epsilon_X}
\circ
{}^g\!\sh{\id_X\odot\bar\epsilon_Y\odot\id_{{}^*X}}
\circ
{}^g\theta_{{}^*Y\odot{}^*X,\,X\odot Y}
\circ
{}^g\!\sh{\id_{{}^*Y}\odot a_X\odot f_Y}
\circ
{}^g\!\sh{\bar\eta_Y}.
\end{align*}

The twisted hexagon axiom, applied to
\(
M={}^*Y\odot{}^*X, N=X, P=Y,
\)
gives
\[
{}^g\theta_{{}^*Y\odot{}^*X,\,X\odot Y}
=
{}^g\theta_{Y\odot{}^*Y\odot{}^*X,\,X}
\circ
{}^g\theta_{{}^*Y\odot{}^*X\odot X,\,Y}.
\]
Naturality of \({}^g\theta\) gives
\begin{align*}
&{}^g\theta_{{}^*Y\odot{}^*X\odot X,\,Y}
\circ
{}^g\!\sh{\id_{{}^*Y}\odot a_X\odot f_Y}
  =
{}^g\!\sh{\id_{Y\odot{}^*Y}\odot a_X}
\circ
{}^g\theta_{{}^*Y,Y}
\circ
{}^g\!\sh{\id_{{}^*Y}\odot f_Y},
\end{align*}
while naturality with respect to \(\bar\epsilon_Y\) gives
\[
{}^g\!\sh{\id_X\odot\bar\epsilon_Y\odot\id_{{}^*X}}
\circ
{}^g\theta_{Y\odot{}^*Y\odot{}^*X,\,X}
=
{}^g\theta_{{}^*X,X}
\circ
{}^g\!\sh{\bar\epsilon_Y\odot
\id_{{}^*X\odot X}}.
\]
Substituting these two identities and using functoriality of the
twisted shadow, we obtain
\begin{align*}
\tr^g(f_X\odot f_Y)
={}&
{}^g\!\sh{\bar\epsilon_X}
\circ
{}^g\theta_{{}^*X,X}
\circ
{}^g\!\sh{a_X}
\circ
{}^g\!\sh{\bar\epsilon_Y}
\circ
{}^g\theta_{{}^*Y,Y}
\circ
{}^g\!\sh{\id_{{}^*Y}\odot f_Y}
\circ
{}^g\!\sh{\bar\eta_Y}.
\end{align*}
Finally,
\(
{}^g\!\sh{a_X}
=
{}^g\!\sh{\id_{{}^*X}\odot f_X}
\circ
{}^g\!\sh{\bar\eta_X}.
\)
Therefore the preceding composite is exactly
\(
\tr^g(f_X)\circ\tr^g(f_Y).
\)
\end{proof}

\end{document}
